%\newpage
%\setcounter{page}{1}
\begin{center}
\begin{figure}
\centering%
\epsfig{file=HojaTitulo/EscudoUN.eps,scale=1}%
\end{figure}
\thispagestyle{empty} \vspace*{2.0cm} \textbf{\huge
Modelo de regresi\'{o}n mixto para datos proporcionales inflados con ceros y/o unos}\\[5.5cm]
\Large\textbf{Juan Camilo D\'{\i}az Zapata}\\[5.5cm]
\small Universidad Nacional de Colombia\\
Facultad de Ciencias, Escuela de Estad\'{\i}stica\\
Medell\'{\i}n, Colombia\\
2017\\
\end{center}

\newpage{\pagestyle{empty}\cleardoublepage}

\newpage
\begin{center}
\thispagestyle{empty} \vspace*{0cm} \textbf{\huge
Modelo de regresi\'{o}n mixto para datos proporcionales inflados con ceros y/o unos}\\[3.0cm]
\Large\textbf{Juan Camilo D\'{\i}az Zapata}\\[3.0cm]
\small Tesis presentada como requisito parcial para optar al
t\'{\i}tulo de:\\
\textbf{Magister en Ciencias Estad\'{\i}sticas}\\[2.5cm]
Director:\\
Ph.D. Freddy Hern\'{a}ndez Barajas\\[2.0cm]
L\'{\i}nea de Investigaci\'{o}n:\\
Modelos de Regresi\'{o}n Mixtos\\[2.5cm]
Universidad Nacional de Colombia\\
Facultad de Ciencias, Escuela de Estad\'{\i}stica\\
Medell\'{\i}n, Colombia\\
2017\\
\end{center}

\newpage{\pagestyle{empty}\cleardoublepage}

\newpage
\thispagestyle{empty} \textbf{}\normalsize
\\\\\\%
\textbf{(Dedicatoria)}\\[4.0cm]

\begin{flushright}
\begin{minipage}{12cm}
    \noindent
        \small
				\textsl{Este trabajo es dedicado a mis padres, Blanca y Guillermo\\
				y a mi futura esposa Zaret.}\\
\end{minipage}
\end{flushright}

\newpage{\pagestyle{empty}\cleardoublepage}

\newpage
\thispagestyle{empty} \textbf{}\normalsize
\\\\\\%
\textbf{\LARGE Agradecimientos}
\addcontentsline{toc}{chapter}{\numberline{}Agradecimientos}\\\\

A Dios, por regalarme salud, tiempo, sabidur\'{\i}a e inteligencia para poder realizar este trabajo, que toda la gloria sea para \'{e}l.\\

A mi tutor, Freddy Hern\'{a}ndez, por el constante apoyo, paciencia, orientaci\'{o}n y confianza depositada a lo largo de estos dos a\~{n}os.\\

A mis padres, Blanca y Guillermo por la paciencia, cari\~{n}o y apoyo d\'{\i}a tras d\'{\i}a.\\

A mi prometida, Zaret, por su paciencia, apoyo incondicional y comprensi\'{o}n; Por ayudarme, escucharme y alentarme en los momentos m\'{a}s dif\'{\i}ciles y as\'{\i} poder ver realizado este trabajo.\\

A mis sobrinos, Jos\'{e} Manuel y Julieta, por la comprensi\'{o}n y el tiempo regalado, a su corta edad respetaron cada uno de los espacios necesarios para elaborar este trabajo.\\

A mi jefe, Andr\'{e}s Ibarra, por haberme brindado los datos para realizar las aplicaciones a datos reales y el tiempo, para poder elaborar este trabajo durante estos dos a\~{n}os.\\

A la universidad de Antioquia, por brindarme sus instalaciones, que me permitieron trabajar por buen tiempo en este trabajo. 

\newpage{\pagestyle{empty}\cleardoublepage}

\newpage
\textbf{\LARGE Resumen}
\addcontentsline{toc}{chapter}{\numberline{}Resumen}\\\\
El modelo de regresi\'{o}n mixto para datos proporcionales inflados con ceros y/o unos, es un modelo de regresi\'{o}n donde la variable respuesta se encuentra definida a partir de una distribuci\'{o}n para datos proporcionales, como la distribuci\'{o}n beta o la distribuci\'{o}n simplex, esta variable puede dar resultados de cero o uno, representando la ausencia o presencia total de cierta caracter\'{\i}stica. Diferentes autores han trabajado en el desarrollo de diferentes modelos y metodolog\'{\i}as de estimaci\'{o}n, sin embargo, no se ha desarrollado un modelo de regresi\'{o}n mixto para datos proporcionales inflados con ceros y/o unos, que re\'{u}na los principales modelos de regresi\'{o}n de este tipo y que la estimaci\'{o}n de los par\'{a}metros sea v\'{\i}a m\'{a}xima verosimilitud y la cuadratura de Gauss-Hermite. En este trabajo se presenta el paquete \pkg{ZOIP} del sistema computacional \proglang{R}, en el que se implementa la distribuci\'{o}n ZOIP (Zeros Ones Inflated Proporcional), el modelo de regresi\'{o}n para efectos fijos y mixtos, ZOIP, que re\'{u}ne las distribuciones y los modelos de regresi\'{o}n de efectos fijos y mixtos para las distribuciones beta y simplex inflada con ceros y/o unos, la estimaci\'{o}n de los par\'{a}metros se hace v\'{\i}a m\'{a}xima verosimilitud y la cuadratura de Gauss-Hermite. Se realizan diferentes estudios de simulaci\'{o}n que muestran la convergencia de los par\'{a}metros y la mejor alternativa de estimaci\'{o}n que presenta mejor desempe\~{n}o. Adem\'{a}s, se presenta el ajuste de diferentes modelos de regresi\'{o}n ZOIP a datos reales.\\

\textbf{\small Palabras clave:} modelos lineales mixtos, datos proporcionales inflados, cuadratura de Gauss-Hermite,  m\'{a}xima verosimilitud.\\


\textbf{\LARGE Abstract}\\\\
El modelo de regresi\'{o}n mixto para datos proporcionales inflados con ceros y/o unos, es un modelo de regresi\'{o}n donde la variable respuesta se encuentra definida a partir de una distribuci\'{o}n para datos proporcionales, como la distribuci\'{o}n beta o la distribuci\'{o}n simplex, esta variable puede dar resultados de cero o uno, representando la ausencia o presencia total de cierta caracter\'{\i}stica. Diferentes autores han trabajado en el desarrollo de diferentes modelos y metodolog\'{\i}as de estimaci\'{o}n, sin embargo, no se ha desarrollado un modelo de regresi\'{o}n mixto para datos proporcionales inflados con ceros y/o unos, que re\'{u}na los principales modelos de regresi\'{o}n de este tipo y que la estimaci\'{o}n de los par\'{a}metros sea v\'{\i}a m\'{a}xima verosimilitud y la cuadratura de Gauss-Hermite. En este trabajo se presenta el paquete \pkg{ZOIP} del sistema computacional \proglang{R}, en el que se implementa la distribuci\'{o}n ZOIP (Zeros Ones Inflated Proporcional), el modelo de regresi\'{o}n para efectos fijos y mixtos, ZOIP, que re\'{u}ne las distribuciones y los modelos de regresi\'{o}n de efectos fijos y mixtos para las distribuciones beta y simplex inflada con ceros y/o unos, la estimaci\'{o}n de los par\'{a}metros se hace v\'{\i}a m\'{a}xima verosimilitud y la cuadratura de Gauss-Hermite. Se realizan diferentes estudios de simulaci\'{o}n que muestran la convergencia de los par\'{a}metros y la mejor alternativa de estimaci\'{o}n que presenta mejor desempe\~{n}o. Adem\'{a}s, se presenta el ajuste de diferentes modelos de regresi\'{o}n ZOIP a datos reales.\\

\textbf{\small Keywords:} linear mixed models, proportional inflated data, Gauss-Hermite quadrature, maximum likelihood.\\